\section{The model}
\label{sec:model}

\subsection{Overview}
The agent is a recurrent vision transformer trained by reinforcement learning to
emit a binary decision (\emph{wait} or \emph{declare change}) on each frame of a
spatially cued orientation-change-detection task (Section~\ref{sec:env}). Three
components are connected in a loop (Figure~\TODO{fig:model}). (1) A patch
front-end parses each frame into a set of spatial tokens. (2) A recurrent
cross-attention encoder integrates the current tokens with two per-token recurrent
memories, maintained by long short-term memory (LSTM) cells across frames
\citep{hochreiter1997long, beck2024xlstm}. (3) An actor--critic reinforcement-
learning module reads the encoder output to select an action and estimate value.
The defining feature of the architecture is that the self-attention is
\emph{recurrent}: attention at each frame is computed not only from the immediate
image but from the contents of a capacity-limited spatial working memory, so that
the cue, shown briefly and then removed, can be maintained and used to bias later
processing \citep{awh2006interactions, panichello2021attvwm, knudsen2007fundamental}.

The encoder is organised into two parallel pathways. The \emph{$\mu$ pathway} is
grounded in the live image and feeds the actor (policy); the \emph{$q$ pathway} is
grounded in accumulated memory and feeds the critic (value). The remainder of this
section gives the precise form of the two pathways and the three structural
configurations we compare.

\subsection{Patch tokens and recurrent memories}
Each $50\times50$ frame $\mathbf{O}^{(t)}$ is divided into a grid of non-overlapping
patches, each embedded by a shared per-patch network and tagged with positional and
temporal codes, yielding a set of tokens
$\Xt^{(t)} = \{\mathbf{x}^{(t)}_i\}_{i=1}^{N}$ with $\mathbf{x}^{(t)}_i \in
\mathbb{R}^{d}$. Two recurrent memories are carried as per-token hidden states (one
slot per patch location), $\Hone^{(t)},\Htwo^{(t)} \in \mathbb{R}^{N\times d}$,
updated frame to frame by independent LSTM cells. $\Hone$ is a fast sensory trace
written from the raw image; $\Htwo$ is a deep trace written from the policy
pathway's own output (below). The two memories play asymmetric roles --- $\Hone$
keeps a recent record of the visual input, $\Htwo$ accumulates the task-relevant
state on which both the decision and its valuation depend.

\subsection{The two cross-attention pathways}
Each pathway is a single-head, pre-norm cross-attention block in which the
\emph{queries} are supplied by the current image tokens and the \emph{keys and
values} are supplied by memory. Writing $\mathrm{Attn}(Q;K)=\mathrm{softmax}
\!\left(QK^{\top}/\sqrt{d}\right)V$ with $V$ tied to $K$, a residual $R$, and a
position-wise feed-forward block $\mathrm{FFN}$, a block computes
\begin{equation}
\mathbf{Z} \;=\; R \;+\; \mathrm{Attn}\big(\Xt;\,\mathrm{KV}\big), \qquad
\mathbf{Z} \;\leftarrow\; \mathbf{Z} + \mathrm{FFN}(\mathbf{Z}).
\label{eq:block}
\end{equation}
The two pathways differ in their key/value bank and residual:
\begin{align}
\text{($\mu$ pathway, $\to$ actor):}\quad
  \Zmu^{(t)} &= \mathrm{block}\big(\Xt^{(t)};\ \mathrm{KV}=[\Hone^{(t)}\!\parallel\!\Htwo^{(t)}],\ R=\Xt^{(t)}\big),
  \label{eq:mu}\\[2pt]
\text{($q$ pathway, $\to$ critic):}\quad
  \Zq^{(t)} &= \mathrm{block}\big(\Xt^{(t)};\ \mathrm{KV}=\Htwo^{(t)},\ R=\Htwo^{(t)}\big).
  \label{eq:q}
\end{align}
The $\mu$ pathway queries \emph{both} memories (concatenated along the token axis,
each source tagged), and carries the raw image on its residual: it is the
grounded, image-reading stream, and its decision can therefore reflect live
visual evidence. The $q$ pathway queries only the deep memory $\Htwo$ and carries
$\Htwo$ on its residual: the image can only \emph{re-weight} the stored state, a
memory-grounded bottleneck appropriate to a contextual value estimate. The memories
are then advanced,
\begin{equation}
\Hone^{(t+1)} = \mathrm{LSTM}_1\!\big(\Xt^{(t)},\Hone^{(t)}\big), \qquad
\Htwo^{(t+1)} = \mathrm{LSTM}_2\!\big(\Zmu^{(t)},\Htwo^{(t)}\big),
\label{eq:mem}
\end{equation}
so that $\Hone$ is written from the raw image and $\Htwo$ is written from the
\emph{policy pathway's} output $\Zmu$. Finally the heads read the two streams with a
\emph{split readout}: the actor reads $\Zmu$ and the (distributional) critic reads
$\Zq$,
\begin{equation}
\pi^{(t)} = \mathrm{Actor}(\Zmu^{(t)}), \qquad
Q^{(t)} = \mathrm{Critic}(\Zq^{(t)}),
\label{eq:heads}
\end{equation}
with the scalar value $V^{(t)}$ derived from the critic's quantile distribution
under the current policy (Section~\ref{sec:methods}).

\subsection{The locus of coupling}
The coupling between policy and value lives in $\Htwo$. The actor's pathway
\emph{writes} $\Htwo$, through $\Zmu$ in Equation~\eqref{eq:mem}, and \emph{reads}
it back, inside the concatenated bank of Equation~\eqref{eq:mu}; the critic's
pathway \emph{reads} the same $\Htwo$ in Equation~\eqref{eq:q}. The two heads are
separated, but they sit on a shared, jointly maintained recurrent memory. This
single shared register --- a common ``workspace'' that the evaluative and the
decision systems both read from and the decision system writes to --- is the
structural variable we manipulate.

\subsection{Three pathway configurations}
\label{sec:configs}
Holding the front-end, the heads, the parameter budget and the entire training
recipe fixed, we compare three wirings (Figure~\TODO{fig:configs}).

\paragraph{Configuration~1 --- shared memory (cross-talk).}
Exactly Equations~\eqref{eq:mu}--\eqref{eq:heads}: two pathways coupled through the
shared $\Htwo$, with the grounded $\mu$ pathway driving the actor.

\paragraph{Configuration~2 --- independent (no cross-talk).}
The same two pathway \emph{types}, but each is given its own private memory, updated
only by its own output:
\begin{align}
\Zmu^{(t)} &= \mathrm{block}\big(\Xt^{(t)};\ \mathrm{KV}=\Hmu^{(t)},\ R=\Xt^{(t)}\big), &
\Hmu^{(t+1)} &= \mathrm{LSTM}_\mu(\Zmu^{(t)},\Hmu^{(t)}),\\
\Zq^{(t)} &= \mathrm{block}\big(\Xt^{(t)};\ \mathrm{KV}=\HQ^{(t)},\ R=\HQ^{(t)}\big), &
\HQ^{(t+1)} &= \mathrm{LSTM}_q(\Zq^{(t)},\HQ^{(t)}).
\end{align}
The modules never touch each other's state; the only thing they share is the raw
input image. This is the minimal edit that removes the cross-talk while preserving
the split, the pathway types, the parameter count and the training recipe.

\paragraph{Configuration~3 --- swapped routing.}
Architecturally identical to Configuration~1, but the readout is reversed: the $q$
(memory) pathway drives the actor and the $\mu$ (image) pathway drives the critic.
This isolates whether it matters \emph{which} pathway commands the policy,
independent of the presence of coupling.

The three configurations differ only in inter-pathway coupling and readout routing.
The comparison therefore isolates the causal contribution of organisation, in the
same spirit as a reversible inactivation experiment
\citep{wardak2006contribution, lovejoy2010inactivation}.
